\subsection{Green's Theorem}

\content{
\begin{theorem} (Green's Theorem) Let $C$ be a positively oriented,
piecewise-smooth, simple closed curve in the plane and let $D$ be the region
bounded by $C$. If $P$ and $Q$ have continuous partial derivatives on an open
region that contains $D$, then 
$$\int_C P\,dx+Q\,dy = \iint_D \left(\pd{Q}{x}-\pd{P}{y}\right)\,dA. $$
\end{theorem}
}{% presentation
}{% nopresentation
}{% comments
}

\content{
\begin{example} Evaluate $\int_C x^4dx+xydy$ where $C$ is the triangular curve
consisting of the line segments from $(0,0)$ to $(1,0)$, from $(1,0)$ to
$(0,1)$, and from $(0,1)$ back to $(0,0)$.
\end{example}
}{% presentation
\vspace{3in}
}{% nopresentation
}{% comments
Answer: $\frac{1}{6}$.
}

\content{
\begin{example} Evaluate $\oint_C (3y-e^{\sin x})dx+(7x+\sqrt{y^4+1})dy$, where
$C$ is the circle $x^2+y^2=9$.
\end{example}
}{% presentation
\vspace{3in}
}{% nopresentation
}{% comments
Mention that the notation $\oint$ is used to denote that the integral is being
taken around a closed, positively oriented curve. 

Also mention that in the book, there is an interesting section on finding areas
with green's theorem, which is the basis of how a planimeter works. 
}

\content{
\begin{example} Evaluate $\int_C y^2\,dx + 3xy\,dy$ where $C$ is the boundary of
the semiannular region $D$ in the upper half plane between the circles
$x^2+y^2=1$ and $x^2+y^2=4$.
\end{example}
}{% presentation
\vspace{3in}
}{% nopresentation
}{% comments
Answer: 14/3
}

